\textbf{IX -- PEZZI FUORI GIOCO}

Talvolta è dato di vedere posizioni in cui un pezzo viene confinato
senza rimedio in una zona delimitata della scacchiera, il possessore di
detto pezzo è in genere destinato alla sconfitta perché gioca le
rimanenti mosse della partita con un pezzo in meno.

In certe occasioni la messa fuori gioco di pezzi avversari può diventare
un obiettivo strategico. Non succede mai che qualcuno si ponga a priori
come piano di ingabbiare un pezzo nemico ma quando un pezzo avversario
ha scarsa mobilità o ne ricorrono comunque gli elementi, tra più mosse
che contribuiscono a mettere in atto il proprio piano, qualsiasi esso
sia, si sceglieranno quelle utili a imbottigliare del tutto il pezzo
avversario.

Winter--Capablanca, Hastings 1919 \textbf{(Partita dei Quattro Cavalli)}

\textbf{1.e4 e5 2.¤f3 ¤c6 3.¤c3 ¤f6} apertura molto popolare all'inizio
del secolo, la Quattro Cavalli è oggi poco giocata. Le si preferisce la
Partita Spagnola, vera regina dei giochi aperti.

\textbf{4.¥b5} il Bianco non ottiene vantaggi a giocare 4.d4 cui può
seguire 4... exd4 5.¤xd4 ¥b4 6.¤xc6 bxc6 7.¥d3 d5 8.exd5 ¤xd5 9.0 0 0-0
l0.¥g5 ¥e6 11.¤b5 c5 12.a3 ¥a5 13.b4 ¥b6 14.bxc5 ¥xc5 con gioco
equivalente (Gheorghiu Suba, Praga 1985).

\textbf{4... ¥b4} l'alternativa è 4... ¤d4 5.¥a4 ¥c5 6.¤xe5 0-0 7.¤d3
¥b6 8.e5 ¤e8 9.0 0 d6 l0.exd6 ¤f6 11.d7 ¥xd7 12.¥xd7 £xd7 13.Cel ¦ae8
14.d3 ¤g4 con parità secondo Estrin.

\textbf{5.0-0 0-0 6.¥xc6 £xc6 7.d3 ¥d6 8.¥g5} il modo corretto di
procedere è ben indicato da Capablanca. Vale la pena di riportare le sue
parole: \emph{``Questa mossa non è affatto in armonia con la natura di
questa variante. Il piano strategico generale del Bianco è di giuocare
h3 e farla seguire, a suo tempo, dalla spinta del pedone g2 in g4 e dal
salto del Cavallo c3 in f5 via e2 e g3, oppure via d1 ed e3. Poi,
possibilmente, il Cavallo f3 sarebbe collegato con l'altro Cavallo
portandolo in h4, h3 o e3, a seconda dell'opportunità. Il Re Bianco
rimane, talvolta, in e1, altra volta vien portato in g2, ma il più delle
volte in h1. Finalmente, il più delle volte, si gioca f4, e quindi
comincia il vero attacco. Talvolta si tratta di un attacco diretto
contro il Re, altra volta si gioca per ottenere un vantaggio di
posizione nel finale, dopo cambiati la maggior parte dei pezzi''.}

\textbf{8... h6 9.¥h4 c5} impedisce 10.d4.

\textbf{10.¤d5? g5 11.¤xf6+ £xf6 12.¥g3} se 12.¤xg5 ¤xd5!

\textbf{12... ¥g4 13.h3 ¥xf3 14.£xf3 £xf3 15.gxf3 f6}

\bookdiagramplaceholder{assets/diagrams/image107.png}{107}

L'¥g3 è imprigionato; la sua liberazione costerebbe almeno un pedone e
molti tempi. In queste condizioni il resto della partita è giocato da
Nero con un pezzo in più.

\textbf{16.¢g2 a5 17.a4 ¢f7 18.¦h1 ¢e6 19.h4 ¦fb8} bloccato il lato di
Re Capablanca comincia le operazioni a ovest.

\textbf{20.hxg5 hxg5 21.b3 c6 22.¦a2 b5 23.¦ha1 c4} tutte le spinte di
rottura sono effettuate. L'apertura del gioco costringerà il Bianco a
una rapida capitolazione.

\textbf{24.axb5 ¤xb3 25.¤xb3 ¦xb5 26.¦a4} il Bianco potrebbe abbandonare
senza rimpianti.

\textbf{26... ¦xb3 27.d4 ¦b5 28.¦c4 ¦b4 29.¦xc6 ¦xd4 30. il Bianco
abbandona.}

Andreoli--Leoncini, Alba Adriatica 1992 \textbf{(Difesa Siciliana)}

\textbf{1.e4 c5 2.d4 ¤xd4 3.c3} nel dicembre del 1989, in un articolo
apparso su \emph{Shakhmatny Bulletin}, due candidati maestri sovietici,
Maliscev e Zipkov, resero nota una variante molto aggressiva contro il
Gambetto Morra: 3... £xc3 4.¤xc3 ¤c6 5.¤f3 e6 6.¥c4 ¤f6 7.0-0 £c7 8.£e2
¤g4! con la minaccia 9... ¤d4! che vince di colpo. In simili occasioni
si è tentati di provare la preparazione dell'avversario col rischio di
ritrovarsi in terreni sconosciuti.

\textbf{3... d5 4.exd5 £xd5 5.¤xd4 ¤c6 6.¤f3 e5 7.¤c3 ¥b4 8.¥d2} una
continuazione molto interessante è 8.¥e2 con l'intenzione di sacrificare
un pezzo dopo 8... e4 (8... exd4!) 9.0-0 ¥xc3 10.bxc3 exf3 11.¥xf3 con
forte iniziativa, per esempio: 11... £a5 12.¦e1+ ¥e6 13.d5 0-0-0 14.£xc6
con vantaggio decisivo.

\textbf{8... ¥xc3 9.¥xc3 e4 10.¤e5 ¤xe5 11.£xe5 ¤e7 12.£c2} una valida
alternativa è 12.£a4+, con l'idea di ostacolare l'arrocco nero. 12...
¥d7 13.£a3 £e6 14.£d4.

\textbf{12... 0-0 13.¦d1 £xa2 14.£xe4 ¥f5 15.¥c4} dopo 15.£xb7 il Bianco
rimane in vantaggio in caso di 15... ¤g6 16.£a6 £xa6 17.¥xa6 ¤xe5 18.¦d5
f6 19.0-0 ¥e6 20.¦a5 ¤c4 21.¦c5, ma il Nero può giocare 15... ¦ad8!
16.£xe7 \emph{(16.¦xd8 ¦xd8 17.£xe7 £d1+ 18.¢e2 £d3+ 19.¢e1 £d1 matto)}
16... ¦xd1+ 17.¢xd1 £d1+ 18.¢e2 £d3+.

\textbf{15... ¥xe4 16.¥xa2 ¦fd8} preferibile alla presa in g2 che
conduce a un gioco piuttosto confuso e di difficile valutazione: 17.¦g1
¥c6 18.e6 f6 (Nun--Witkowski, Hradec Kralove 1975).

\textbf{17.0-0} con questa mossa entrambi i giocatori hanno completato
lo sviluppo. In questi casi, prima di muovere, è sempre bene fermarsi un
attimo e fare qualche considerazione di ordine generale. Le schermaglie
d'apertura non sono state del tutto soddisfacenti per il Nero perché
hanno lasciato all'avversario la coppia degli Alfieri; poiché la
posizione è piuttosto aperta, in caso di gioco superficiale questo
elemento può diventare decisivo. Tutti e due i giocatori hanno, inoltre,
una Torre sulla colonna ``d'', questo significa che chi per primo cambia
la Torre, pe¢de la colonna. Per queste ragioni il Nero gioca:

\textbf{17... ¥d5} se il Bianco cambia il pezzo pe¢de il vantaggio della
coppia degli Alfieri; se sposta l'Alfiere per evitare il baratto si può
giocare l'Alfiere in b3 e conquistare la colonna.

\textbf{18.¥b1} decidendo di tenersi gli Alfieri \ldots{}

\textbf{18... ¥b3 19.¦xd8+ ¦xd8}\ldots{} ma cedendo la colonna.

\textbf{20.¥b4} il Bianco contava su questa mossa per chiudere la
colonna ma non è così facile.

\textbf{20... ¤c6 21.¥d6} con la sicurezza con cui eseguì questa mossa
Mark Andreoli sembrava molto contento di sé.

\textbf{21... ¤d4}

\bookdiagramplaceholder{assets/diagrams/image108.png}{108}

Non tanto per centralizzare il Cavallo quanto per andare in b5 e
attaccare l'¥d6. Il Nero migliora l'assetto dei propri pezzi senza
dimenticarsi che non sarebbe male togliere all'avversario uno dei due
Alfieri.

\textbf{22.¥e4 b6 23.¦a1} meccanica, non ostacola il piano del Nero che
è di andare di Cavallo in b5. In partita ebbi l'impressione che il mio
avversario non si aspettasse che rimuovessi il Cavallo, così ben
centralizzato.

\textbf{23... ¤b5 24.f3} stiamo per entrare in un finale dove le Torri
avranno un ruolo preminente; il Bianco apre una finestrella per il Re.

\textbf{24... ¤xd6 25.exd6 ¦xd6 26.¦xa7} raggiunti gli obiettivi occorre
riesaminare la posizione e formulare un piano. Il Bianco non vanta più
alcun vantaggio, anzi, il Pb2 sembra più debole del Pb6. Prossimo
obiettivo sarà farlo cadere. Se prima si giocava per annullare il
vantaggio del Bianco ora si comincia a giocare per vincere!

\textbf{26... g6 27.¦b7 f5} utile a scoordinare un po' i pezzi bianchi.

\textbf{28.¦b8+ ¢g7 29.¥a8 ¥c4} preparandosi a sostenere il pedone
\emph{b} con l'Alfiere, manovra necessaria per liberare la Torre che
deve andare a catturare il Pb2.

\textbf{30.¦b7+ ¢f6} non curandosi del Ph7. La difesa del Ph7 non vale
la messa fuori gioco del Re (30... ¢h6 31.f4).

\textbf{31.¦c7 ¦d1+ 32.¢f2 b5} ora si può finalmente procedere alla
cattura del Pb2.

\textbf{33.f4} impedisce di giocare f4 al Nero.

\textbf{33... ¦d2+ 34.¢e3 ¦d3+ 35.¢f2 ¦b3 36.¦xh7 ¦xb2+} un altro
importante obiettivo è raggiunto. Ora comincia l'ultimo atto: promuovere
il Pb5.

\textbf{37.¢e3 ¦e2+ 38.¢d4 ¦f2 39.g3 ¥f7 40.¢c3 ¦e2 41.¢b4 ¦b2+ 42.¢c3
¦b3+ 43.¢c2 ¦a3} c'è sempre modo di pe¢dere: 43... b4?? 44.¦xf7+!

\textbf{44.¥c6 b4 45.¥b5 ¦a2+ 46.¢b1 ¦a8} prima di eseguire ulteriori
spinte il Nero prova a tagliare definitivamente fuori gioco la Torre
bianca.

\textbf{47.¢b2 ¥g8 48.¦h8 ¢g7 49.¦h4?} pur di non far cadere un pedone
il Bianco perde l'ultima occasione di far rientrare in gioco la Torre:
49.¦b8 ¦a2+ e il Ph2 sarebbe caduto. L'avarizia sarà presto punita.

\textbf{49... ¦a2+ 50.¢b1 ¦f2!}

\bookdiagramplaceholder{assets/diagrams/image109.png}{109}

Frustrando
le speranze del Bianco di rimettere in gioco la Torre tramite la spinta
in g4, cui ora seguirebbe ¦xf4. Ora il Bianco gioca con una Torre in
meno e non può evitare la sconfitta.

\textbf{51.¥a4 ¥c4 52.¥c6 ¥d3+ 53.¢a1 b3} il Bianco abbandona; non c'è
modo di evitare b2 matto.

\bookdiagramplaceholder{assets/diagrams/image110.png}{110}

La manovra per vincere non è difficile da trovare: £d2, ¢b3, ¤b2, ¢a4 e
il Pa5 cade.

Il Nero giocò:

\textbf{37... £d6!} mettendo in presa addirittura il pezzo più grosso.

\textbf{38.¤xb6? ¤xb6 39.h4} il Bianco si rende conto del rischio di una
chiusura della posizione con 38... h4 e conta su questa mossa per
entrare con la Donna nello schieramento avversario via h3, ma il Nero
aveva visto più lontano.

\textbf{39... gxh4 40.£d2} con l'idea £g2-h3.

\textbf{40... h3!} completamente inaspettata! La minaccia 41... h2
impone la presa.

\textbf{41.gxh3 h4!} e il Nero ha confinato la Donna avversaria in metà
scacchiera: patta!

Per finire vale la pena di dare un'occhiata a questo studio, paradossale
ma divertente.

\bookdiagramplaceholder{assets/diagrams/image111.png}{111}

\emph{Il Bianco muove e patta}

\textbf{1.¥a4+! ¢xa4} c'è lo scacco perpetuo in caso di 1... ¢c4 2.¥b3+
¢b5 3.¥a4+ ¢c4 4.¥b3+.

\textbf{2.b3+ ¢b5 3.c4+ ¢c6 4.d5+ ¢d7 5.e6+!} si priva dell'ultimo
pezzo.

\textbf{5... ¢xd8 6.f5} chiude in modo definitivo la porta e rende
inutili i tre pezzi del Nero.
